\chapter{Vở Sạch, Vở Ghi, Vở Trắng}
\chapterintro{Có những cuốn vở sạch đến mức khiến người ta nghi ngờ cả một học kỳ đã đi đâu}{Vở trắng đôi khi không phải vì nghèo chữ, mà vì quá giàu trì hoãn.}

\begin{lucbatbox}{Lục bát: Vở mới quanh năm}
Đầu năm em bọc rất xinh\
Giữa năm nhìn lại vẫn tinh như đầu\
Trang nào cũng trắng một màu\
Như chưa từng phải đi qua tháng ngày\
Bạn nhìn mà dạ đắng cay\
Vở em giữ đẹp bằng tay khỏi dùng
\end{lucbatbox}

\begin{tutuyetbox}{Thất ngôn tứ tuyệt: Chữ hiếm như vàng}
Vở em quý chữ vô ngần lắm
Nên từng con chữ giữ như không
Ai nhìn cũng thấy quá mênh mông
Một trời giấy trắng khiến ngẩn ngơ
\end{tutuyetbox}

\begin{ghichu}
Vở trắng sạch đẹp thật, nhưng trong bối cảnh đi học lâu ngày thì vẻ đẹp ấy thường hơi đáng nghi.
\end{ghichu}

\ngancach

\begin{lucbatbox}{Lục bát: Chỉ ghi tiêu đề}
Tiêu đề em viết rất xinh\
Mực màu, nét thẳng, khung hình chỉn chu\
Qua phần nội dung tới lui\
Trang vở bỗng rộng như trời tháng ba\
Người xem cũng phải xuýt xoa\
Đẹp thì thật đẹp, chỉ là hơi trống
\end{lucbatbox}

\begin{tutuyetbox}{Thất ngôn tứ tuyệt: Vở đẹp mà buồn}
Vở em sạch đến nao lòng ghê
Như chưa đi hết dòng giảng nào
Nếu thêm chữ nghĩa vài ba dòng
Trang giấy ấy hẳn bớt sầu hơn
\end{tutuyetbox}